\setcounter{section}{128}
\section{Continuidad}

Nombre completo: Temas Específicos para la preparación de la Oposición al Cuerpo Superior de Sistemas y Tecnologías de la Información de la Administración del Estado.

Origen: ASTIC

\localtableofcontents

\subsection{Gestión de la continuidad del negocio} %3
    La norma ISO 22301 \textit{Security and Resilience - Business continuity management systems - Requirements} especifica los requisitos de establecimiento y gestión de un  SGCN (sistema de gestión de la continuidad de negocio) eficaz.

    Todo SGCN se debe basar en:
    \begin{itemize}
        \item análisis de las necesidades de la organización y beneficio de la política y objetivos de la continuidad del negocio
        \item Implantación y ejecución de controles y medidas para gestionar la capacidad de la organización frente a incidentes que amenacen la continuidad del servicio
        \item Supervisión y revisión del rendimiento del SGCN
        \item Mejora contínua
    \end{itemize}

    Todo SGCN posee estos componentes principales. Todos los procesos de SGCN son importantes en la organización:
    \begin{itemize}
        \item Política
        \item Personas con responsabilidad bien delimitada
        \item Procesos de gestión de \begin{itemize}
            \item Política
            \item Planificación
            \item Implantación y operación
            \item Evaluación del rendimiento
            \item revisión por la dirección
            \item mejora
        \end{itemize}
        \item Pruebas auditables documentadas
    \end{itemize}

    Los SGCN se basan en el modelo Planificar - Hacer - Verificar - Actuar. (Plan, do, check, act):
    \begin{description}
        \item[Planificar] definir la política de continuidad de negocio, los objetivos y metas, controles, procesos y procedimientos para mejorar la continuidad y obtener resultados alineados con las políticas y objetivos de la organización
        \item[Hacer] (Implantar y operar) la política, controles, procesos y procedimientos de la continuidad del negocio.
        \item[Verificar]Supervisar el rendimiento en base a objetivos y política de continuidad de negocio. Implica también informar de los resultados y definir medidas para corrección y mejora
        \item[Actuar] Mantener y mejorar el SGCN (bucle de feedback)
    \end{description}


    \subsubsection{Contexto de la organización (Planificar)} %5
        Se trata de establecer el contexto del SGCN, con sus necesidades, requisitos y alcance.

        En cuanto a la operativa de la organización se debe identificar y documentar:
        \begin{itemize}
            \item Actividades, funciones, servicios, productos, asociaciones, flujos de trabajo, relaciones con stakeholders y su impacto ante un incidente
            \item Las relaciones entre la política de continuidad de negocio y los objetivos de la organización, políticas y estrategia global de gestión de riesgos.
        \end{itemize}

        Todo contexto debe:
        \begin{itemize}
            \item Definir objetivos de forma clara, incluidos los relativos a la continuidad del negocio
            \item Establecer factores y fuentes internas y externas que provocan riesgos
            \item Calcular la probabilidad e impacto de los riesgos
            \item Definir la finalidad del SGCN \begin{itemize}
                \item Qué partes de la organización incluir
                \item requisitos del SGCN según misión, metas, obligaciones, responsabilidades legales y reglamentarias
                \item Identificar productos, servicios y actividades relacionados con el SGCN
                \item necesidades de los stakeholders
                \item tener en cuenta el tamaño, naturaleza y complejidad de la organización
            \end{itemize}
        \end{itemize}

    \subsubsection{Liderazgo (Planificar)} % 5
        La alta dirección debe estar comprometida y liderar el SGCN. Esto debe plasmar:
        \begin{itemize}
            \item posibilitar políticas y objetivos del SGCN según la estrategia de la organización
            \item Fomentar la integración de requisitos del SGCN en los procesos de negocio
            \item Dotar de recursos al SGCN
            \item Que la continuidad de negocio sea permeable en la organización, fomentar que los mandos intermedios demuestren liderazgo y compromiso.
            \item Realizar seguimiento de los resultados del SGCN
            \item Dirigir y apoyar a las personas encargadas del SGCN
            \item Fomentar la mejora contínua
        \end{itemize}

        Para ello cuenta con las siguientes herramientas y procesos para cubrir lo anterior, y facilitar el establecimiento, operación, supervisión mantenimiento y mejora del SGCN:
        \begin{itemize}
            \item Política de Continuidad de negocio
            \item Objetivos y planes del SGCN
            \item Definir y normalizar puestos con funciones, responsabilidades y competencia de la gestión de continuidad de negocio. Asegurando la formación contínua.
            \item Determinar responsables del SGCN con autoridad y funciones para implantar y mantener el SGCN
        \end{itemize}

        La política debe alinearse con los objetivos de la organización, definir la estructura que establezca los objetivos de la continuidad de negocio e incluir la obligatoriedad de requisitos del SGCN y de mejora contínua. Esta política debe estar documentada y disponible. Debe difundirse dentro de la organización y revisarse periódica y sistemáticamente.

        Debe garantizar la asignación de la autoridad y responsabilidades del SGCN.

    \subsubsection{Los objetivos del SGCN (Planificar)} %7
        La dirección debe garantizar que los objetivos del SGCN se establecen y comunican en la empresa.

        Cada objetivo debe tener: un responsable, acciones, recursos, planificación y evaluación de resultados con sus indicadores.

    \subsubsection{Competencias, Comunicación y Documentación: Apoyo (Planificar)} %7
        La planificación debe describir las operaciones para definir competencias, comunicaciones con stakeholders y el ciclo de vida de la documentación del SGCN.

        Las competencias deben:
        \begin{itemize}
            \item asignar y definir competencias a su personal
            \item Posibilitar la formación inicial, contínua y desarrollo de carrera profesional
            \item Evaluar la formación contínua
            \item Almacenar la documentación que acredite las competencias del personal.
        \end{itemize}

        En cuanto a Concienciación para todos los trabajadores:
        \begin{itemize}
            \item Política de continuidad de negocio
            \item Beneficios y eficacia del SGCN
            \item Consecuencias de no cumplir con los requisitos
            \item Protocolos de actuación ante un incidente
        \end{itemize}

        En cuanto a las comunicaciones, se deben detallar:
        \begin{itemize}
            \item temática de las comunicaciones
            \item cuándo se comunicarán
            \item cómo comunicarse (los medios)
            \item a quien comunicar
            \item quien comunicará
        \end{itemize}

        La organización contará con procedimientos de comunicación para:
        \begin{itemize}
            \item comunicaciones internas y externas
            \item recepción, almacenado y respuesta con stakeholders
            \item Integración de sistema de aviso de amenazas a la continuidad de negocio
            \item Disponibilidad de medios de comunicación ante incidentes y pruebas periódicas de estos
            \item Asegurar la interoperabilidad con otras organizaciones y usuarios.
        \end{itemize}

        Por ello es fundamental el ciclo de vida de la documentación del SGCN. Esta debe incluir la identificación y descripción, unos metadatos que la cataloguen. Y un formato definido.

        Toda documentación debe ser catalogada para asegurar su disponibilidad, confidencialidad, integridad y trazabilidad de uso. Para ello se deben realizar las siguientes operaciones:
        \begin{itemize}
            \item Distribución, acceso, recuperación y utilización
            \item Almacenaje y conservación
            \item Control de versiones y modificaciones
            \item Retención y disposición
            \item Recuperación y utilización
            \item No utilización de información obsoleta
        \end{itemize}

    \subsubsection{Operación (Hacer)} %8
        Se deben planificar, implantar y controlar procesos. Para ello se debe realizar un análisis de impacto y de riesgos.

        %Gif del negro tirando las gafas. Por qué cojones tengo cuatro putas fases de planificar y lo primero que hay en la fase de hacer es planificar mas??????????
        % Porque en las partes de planificar no se planifica, se planifica como planificar y qué tener en cuenta a la hora de planificar

        El análisis de impacto identifica las actividades críticas de la organización, evalúa el impacto de la no realización y define plazos aceptables. El análisis de riesgos identifica los riesgos posibles sobre dichas actividades y la posibilidad de que ocurran.

        En base a estos análisis se define la estrategia de continuidad de negocio con los procedimientos, equipos y planes de continuidad de negocio.

        El análisis de impacto sobre el negocio es necesario para identificar las operaciones y servicios críticos. Debe contener: actividades, evaluación del impacto temporal de la no realización de cada actividad, definición de plazos de reanudación de actividades aceptables y identificación de dependencias y recursos de las actividades.

        El análisis de riesgos debe: identificar los riesgos de las actividades de la organización, analizarlos periódicamente, evaluar los riesgos relacionados con interrupciones e identificar los procedimientos a aplicar para cada riesgo, conforme al SGCN.

        \paragraph{Estrategia de continuidad de negocio}
            La estrategia de continuidad de negocio establece los objetivos y prioridades de recuperación en función de los análisis anteriores. 

            Los objetivos de la estrategia son:
            \begin{itemize}
                \item Salvaguardar las actividades mas importantes
                \item EStabilizar, continuar, reanudar y recuperar dichas actividades, dependencias y recursos
                \item Mitigar, gestionar y responder a los impactos
                \item incluir plazos de reanudación de actividades
                \item Incluir cómo evaluar a los proveedores externos
            \end{itemize}

            Para la implantación de la estrategia, los recursos deben incluir:
            \begin{itemize}
                \item Personas
                \item Informaciones y datos
                \item Edificios, entornos de trabajo y servicios asociados
                \item Instalaciones, equipos y consumibles
                \item Sistemas TIC
                \item Transportes
                \item Finanzas
                \item Socios y proveedores
            \end{itemize}

            Debe incluir medidas proactivas ante los riesgos que requieran tratamientos que disminuyan la probabilidad de corte de servicio o del periodo de interrupción y reduzcan el impacto de corte de servicio sobre los servicios y productos de la organización.

            Definida la estrategia se deben definir los procedimientos y equipos de continuidad de negocio.

        \paragraph{Procedimientos de Continuidad de Negocio}
            los Procedimientos de continuidad de negocio deben establecerse, implantarse y mantenerse para gestionar los incidentes y continuar con las actividades. Deben estar basados en los objetivos de recuperación del análisis de impacto de negocio

            Los procedimientos deben:
            \begin{itemize}
                \item tener un protocolo de comunicaciones internas y externas con asignación de roles y responsabilidades
                \item Contener medidas específicas a realizar durante la interrupción que minimicen las consecuencias
                \item Ser flexibles con las amenazas imprevistas o cambios de condiciones internas o externas
                \item Estar dirigidos a los sucesos que interrumpen operaciones.
            \end{itemize}

            Se deben gestionar los procesos de aviso y comunicación para:
            \begin{itemize}
                \item detectar y supervisar incidentes
                \item Comunicar interna y externamente de forma adecuada
                \item garantizar los sistemas de comunicación durante incidentes y su interoperabilidad
                \item registrar información del ciclo de vida del incidente
                \item probarse periódicamente y garantizar el funcionamiento
            \end{itemize}

            La organización debe someter a pruebas y ensayos sus procedimientos de continuidad de negocio. Estas pruebas deben ser coherentes con el campo de aplicación y objetivos del SGCN. Se deben definir metas y objetivos. Se debe minimizar el riesgo de interrupción. Se debe informar de resultados, recomendaciones y acciones y deben ser revisadas en el contexto de mejora contínua.

        \paragraph{Equipos de continuidad de negocio}
            Se deben gestionar por la organización para responder a incidentes, utilizando personal con autoridad, responsabilidad y competencia. 

            Sus tareas incluyen:
            \begin{itemize}
                \item Evaluar la naturaleza de incidentes e impactos potenciales
                \item Definir umbrales de impacto que justifiquen la respuesta y activación de procedimientos
                \item Activar la respuesta de continuidad de negocio adecuadas
                \item Planificar las acciones que deben ser tomadas
                \item Monitorizar los efectos de incidentes y las respuestas
                \item comunicar las partes interesadas
            \end{itemize}

        \paragraph{Planes de continuidad de negocio}
            Guías detalladas para la implantación de la estrategia de recuperación. Establecen y asignan las funciones, responsabilidades y delegación de poderes en caso de que un suceso de gravedad extrema afecte al personal clave de una entidad.

            La organización debe establecer planes que deben contener:
            \begin{itemize}
                \item Funciones y responsabilidades de los equipos con autoridad ante incidentes
                \item Un proceso que active la respuesta ante incidentes
                \item Cómo gestionar incidentes teniendo en cuenta el bienestar de personas, estrategia de la organización y actividades prioritarias de esta
                \item cómo comunicar los incidentes a empleados, familiares, partes interesadas y servicios de emergencia
                \item Cómo continuar el servicio en tiempo establecido
                \item Cómo dar respuesta del incidente a los medios de comunicación, estrategia de comunicación, medios, directrices, portavoz\dots
                \item proceso de finalización una vez se ha mitigado
            \end{itemize}

            Cada plan debe tener a su vez:
            \begin{itemize}
                \item Finalidad y campo de aplicación
                \item Objetivos
                \item Activación (criterios y procedimiento)
                \item Procedimiento de continuidad que aplican al plan
                \item Funciones, responsabilidades y autoridades
                \item Requisitos y procedimientos de comunicación
                \item Interdependencias y relaciones internas y externas
                \item Requisitos de los recursos
                \item Procesos de flujo de información y comunicación
            \end{itemize}

    \subsubsection{Evaluación del rendimiento (Verificar)} %12
        Para evaluar el rendimiento y eficacia del SGCN se debe definir: qué, cómo cuando supervisar y medir, cómo analizar y evaluar y cuando analizar los resultados.

        Esta evaluación del rendimiento permite prevenir interrupciones. Todas las evaluaciones deben conservarse como evidencias.

        \paragraph{Procedimiento de supervisión}
            Deben cumplir:
            \begin{itemize}
                \item Definir los indicadores de rendimiento
                \item Supervisar el cumplimiento de la política de continuidad, objetivos y metas del SGCN
                \item Supervisar los procedimientos de continuidad de negocio de las actividades prioritarias
                \item Supervisar el cumplimiento de la normativa aplicable
                \item Registrar los resultados de evaluaciones
            \end{itemize}

            De esta forma, la organización debe evaluar los procedimientos de continuidad del negocio de forma periódica. Estos deben ser adecuados y eficaces y conformes a normativas, buenas prácticas, política y objetivos de negocio.

            En cualquier caso, tras un incidente que produce la activación de un procedimiento, se debe hacer una revisión \textit{a posteriori}.

        \paragraph{Auditoría interna}
            Se usan para validar que el SGCN es conforme a requisitos y normativa aplicable, que está implantado y mantenido de forma eficaz.

            La organización debe:
            \begin{itemize}
                \item gestionar programas de auditoría
                \item definir criterios y alcance de las auditorías
                \item Seleccionar auditores objetivos e imparciales
                \item posibilitar que la dirección sea informada de los resultados
                \item conservar los resultados
            \end{itemize}

            Las propuestas de las auditorías deben ser ejecutadas.

        \paragraph{Revisión de la dirección}
            La dirección debe revisar periódicamente el SGCN. Deben tener en cuenta:
            \begin{itemize}
                \item estado de las acciones a llevar a cabo en reuniones anteriores
                \item cambios internos y externos importantes para el SGCN en política y objetivos
                \item Información del rendimiento del SGCN \begin{itemize}
                    \item No conformidades
                    \item acciones correctoras
                    \item resultados de la evaluación de la supervisión y medición
                    \item resultados de auditorías
                \end{itemize}
                \item oportunidades de mejora contínua
                \item Riesgos no cubiertos en el análisis de riesgos anterior
                \item Buenas prácticas y directrices que aparecen
                \item Modificar el SGCN para mejorar su eficacia: \begin{itemize}
                    \item Cambiar objetivos
                    \item Actualización del análisis de riesgos, de impacto, planes de continuidad y procedimientos de continuidad
                    \item Modificación de los procedimientos de continuidad de negocio y control
                \end{itemize}
            \end{itemize}

    \subsubsection{Mejora (Actuar)} %13
        Ante una no conformidad o incidente se debe:
        \begin{itemize}
            \item identificar y reaccionar ante ella, controlando, corrigiendo y afrontando las consecuencias
            \item evaluar la necesidad de aplicar acciones para que no se repita, revisando el incidente, causas y riesgos similares
            \item Aplicar las acciones necesarias y revisar su eficacia
            \item Revisar cambios en el SGCN si fuese necesario
            \item Conservar la información del incidente, naturaleza de este y acciones tomadas
        \end{itemize}

        Además debe revisar y mejorar de manera contínua la idoneidad, adecuación y eficacia del SGCN.

\subsection{Planes de continuidad y contingencia} %15
    \subsubsection{Tipos de planes de continuidad} %15
        Existen tres tipos de planes dentro de los Planes de Continuidad de Negocio. Todos ellos con fases idénticas y cada uno de los señalados contenido en el anterior (PCN con un PCN-TIC con un DRP):
        \begin{description}
            \item[Plan de Continuidad de Negocio] Plan de continuidad ante diversas perspectivas como Infraestructuras TIC, recursos humanos, mobiliario, comunicación, logística, sistemas industriales, etc\dots Cada uno con un plan específico
            \item[Plan de Continuidad Tic] plan referido únicamente a las TIC
            \item[Plan de Recupeeración de Desastres (DRP)] Plan reactivo ante una posible catástrofe.
        \end{description}

    \subsubsection{Fase 0: Determinación del alcance} %16
        Para determinar el alcance y asignar recursos y tiempos óptimos se debe empezar por los departamentos o áreas mas críticas. Siempre con el compromiso de la dirección de ampliar la continuidad a toda la organización. Eventualmente se debe abarcar todo, pero se comienza por las partes prioritarias.

        La fase es la mas corta y la que menos recursos tiene. Determina la magnitud del proyecto y su viabilidad futura. Se estudia qué elementos de la organización son el foco de la mejora contínua. Afecta al personal, activos, sistemas informáticos, servicios y procesos. Tiene alcance a los sistemas o procesos mas críticos, aquellos cuya pérdida generaría mas impacto en la organización. Requiere la colaboración de todos los departamentos implicados.

        El plan puede enfocarse por activos o por procesos. Enfocándolo por activos se selecciona un conjunto de activos y con ellos los procesos que los utilizan. En el enfoque por procesos, se realiza independientemente de los activos que participan.

    \subsubsection{Fase 1: Análisis de la organización} %16
        Se identifican los procesos críticos de negocio, sus activos y sus necesidades temporales y recursos.

        Para ello se deben mantener reuniones con los usuarios finales del proceso para obtener dependencias de proveedores, personal y aplicativo. Se debe recopilar la información sobre las aplicaciones del proceso (detalles del funcionamiento, instalación, proveedor, mantenimiento, copias de seguridad, acuerdos de nivel de servicio).

        Con todo ello se elaboran análisis de impacto sobre negocio y análisis de riesgo

        \paragraph{Análisis de Impacto sobre negocio (BIA)}
            Debe contener los requisitos temporales y de recursos de los procesos dentro del alcance.
            \begin{description}
                \item[Tiempo de ejecución RTO] el tiempo que un proceso estará detenido antes de res restaurado, es algo subjetivo
                \item[RRHH y tecnológicos del proceso] Personal crítico sin reemplazo y aplicaciones, sistemas, equipamiento y elementos auxiliares necesarios para que un proceso funcione en contingencia.
                \item[Tiempo máximo tolerable de caida (MTD)] Tiempo aceptable sin consecuencias desastrosas para la organización. Es subjetivo. MTD es un parámetro de negocio y RTO es técnico. Se debe buscar que RTO < MTD
                \item[Niveles mínimos de recuperación de servicio ROL] nivel mínimo de recuperación de una actividad aunque no opere a niveles óptimos. Puede establecerse en niveles absolutos o porcentuales
                \item[Dependencias de otros procesos internos o proveedores externos] Según la criticidad de las actividades, se puede solicitar los planes de recuperación a los proveedores y los tiempos de recuperación que maneja.
                \item[Grado de dependencia de la actualidad de los datos RPO] Determina el impacto que tiene sobre la actividad la pérdida de datos. Esto influye en las políticas de backups de la empresa. Esta criticidad puede determinarse por procesos.
            \end{description}

            Los vectores son muy subjetivos y se deben contrastar con datos objetivos como el histórico de incidencias, de bajas\dots Estos valores determinan las medidas a tomar y su coste.

        \paragraph{Análisis de riesgos}
            Se deben estudiar las amenazas de los procesos críticos, probabilidades, impactos y activos. Los procesos críticos podrían ser los de un MTD inferior a 24 horas.

            Todo análisis de riesgos pasa por 3 fases
            \begin{description}
                \item[Determinar amenazas] que provocan una indisponibilidad.
                \item[Probabilidad de impacto] se suelen usar escalas subjetivas.
                \item[Producto del impacto por la probabilidad] Se caracteriza el nivel de riesgo según el producto de estas dos medidas.
            \end{description}
            Existen metodologías mas complejas. Pero aquellas que cuentan con una mayor probabilidad e impacto son las que es mas posible que pase y que mas afectan, por lo tanto las primeras a las que atender.

            Todo riesgo debe tratarse de alguna manera (y establecerse en el plan de gestión de riesgos). Pueden transferirse a un tercero eliminarse, asumirlo, o mitigarlo.

    \subsubsection{Fase 2: Determionación de la Estrategia de Continuidad} %18

        Tras las fases 0 y 1 tenemos los procesos críticos, tiempos de recuperación, requisitos de pérdida de datos, recursos, tiempos de recuperación de recursos y riesgos.

        Se debe determinar si es posible restaurar los procesos críticos de negocio en tiempo con los recursos y estrategias actuales. Si no se puede, se deben establecer estrategias de recuperación.

        Los elementos que suelen verse afectados por contingencias son:
        \begin{description}
            \item[Personal] mitigar la ausencia de personal crítico
            \item[Locales]evaluarse situaciones sin lugar para trabajar
            \item[Tecnología] valorar alternativas o medidas complementarias
            \item[Información] disponibilidad y confidencialidad de la información de procesos críticos
            \item[Proveedores] Revisar los acuerdos de nivel de servicio
        \end{description}

        Como resultado habrá distintas estrategias de recuperación por proceso. Deben valorarse según el coste, viabilidad de implantación, mantenimiento, recursos\dots para obtener un conjunto de estrategias a implantar.

    \subsubsection{Fase 3: Respuesta a la contingencia} %18
        Se debe implantar la estrategia. Se comienza con las estrategias para mejorar la continuidad del negocio, clasificación y priorización en base al proceso y su criticidad.

        Se aborda la documentación de respuesta a la contingencia. Se ejecuta en forma de árbol jerárquico, donde la raíz gestiona el momento inmediatamente posterior a la crisis. Cada rama detalla las bases de la recuperación de infraestructura, y las hojas los procedimientos técnicos detallados.

        \paragraph{Plan de crisis}
            Documento con todos los elementos cruciales para la gestión de los momentos iniciales de la crisis. El objetivo es no tomar decisiones improvisadas y equivocadas.

            \begin{description}
                \item[Condiciones de activación]
                \item[Flujos de toma de decisiones]
                \item[Medios para la declaración de la crisis]
                \item[Personal encargado de activar el plan de crisis y su gestión]
                \item[Datos de contacto del personal implicado]
                \item[Niveles de priorización en la recuperación de infraestructura]
                \item[Requisitos temporales de puesta en marcha]
                \item[Planes operativos existentes y personal responsable]
            \end{description}

        \paragraph{Planes Operativos de Recuperación de Entornos}
            Una vez pasado el momento inicial se debe evaluar su alcance y ver qué Planes Objetivos de Recuperación de Entornos se activan. Estos planes contienen uno o mas entornos independientes y su información detallada (Como Correo, ERP, gestor documental).

            Cada plan ordena qué Procedimientos Técnicos de Trabajo deben ser ejecutados.

        \paragraph{Procedimientos Técnicos de Trabajo}
            Es la documentación que indica cómo gestionar y recuperar un sistema. Forman parte de la mejora contínua y del día a día, pero son cruciales en una gestión de crisis. Incluyen versiones de los programas, comandos, recuperación de bases de datos, puesta en marcha y otra información de uso.

    \subsubsection{Fase 4: Prueba mantenimiento y revisión} %19
        Todo plan debe mantenerse actualizado y comprobado. Se consideran buenas prácticas hacerlo anualmente para cubrir todas las catástrofes.
        
        Debe incluirse en las pruebas a los usuarios, personal externo, descripción de la prueba a realizar y su resultado. Fecha y hora de realización intentando impactar lo mínimo posible.

        El resultado de la prueba debe documentarse con las incidencias surgidas. Todas deben ser analizadas y corregidas como proceso de mejora. 

        Algunas pruebas que deben ser realizadas son la caída del suministro eléctrico, tiempos de recuperación de repositorios documentales y permisos. Recuperación de aplicaciones críticas, acceso remoto o disponibilidad de entornos replicados.

        \paragraph{Plan de Mantenimiento}
            Es un elemento en el que mantener actualizada la documentación ante cualquier cambio de la organización. En caso de crisis esta documentación debe estar actualizada.

            Define qué eventos provocan una revisión no planificada (como cambios de proveedor o de un sistema)
        
        \paragraph{Plan de pruebas}
            Contiene las pruebas de contingencia a llevar a cabo en una organización.

            Estos ayudan a mantener información actualizada, continuidad de la organización en tiempo, mejor relación y rendimiento del personal implicado, conocimiento de los usuarios de las pruebas y mayor confianza de los usuarios.

    \subsubsection{Fase 5: Concienciación} %20
        Se debe plantear un proceso de concienciación que contemple la descripción de los elementos que se utilizan en la continuidad.
        \begin{itemize}
            \item análisis de impacto sobre el negocio
            \item plan de crisis
            \item estrategias de recuperación
            \item responsabilidades
            \item plan de pruebas
            \item todo lo relacionado con el plan de contingencia que ayude a mejorar periodos de recuperación
        \end{itemize}

        El público objetivo es tanto el personal técnico como de negocio con relación con los procesos críticos.

% \subsection{Conclusión}v %21